\subsection{Indecomposable modules. Elementary divisors}

DEFINITION 3. --- A left module M over a ring A is said to be decomposable if it is the direct sum of a family of proper nonzero submodules. Otherwise it is said to be indecomposable.

The zero module is thus decomposable, being the direct sum of the empty family of submodules.

Let \(\mathfrak{a}\) be a left ideal of the ring \(\mathbf{A}\) ; the submodules of \(\mathbf{A} / \mathfrak{a}\) are just the quotients \(\mathfrak{b} / \mathfrak{a}\) , where \(\mathfrak{b}\) is an ideal of \(\mathbf{A}\) containing \(\mathfrak{a}\) (I, p.~41, Th. 4); if \(\mathfrak{b}\) and \(\mathfrak{c}\) are two ideals of \(\mathbf{A}\) containing \(\mathfrak{a}\) , then the module \(\mathbf{A} / \mathfrak{a}\) is the direct sum of the submodules \(\mathfrak{b} / \mathfrak{a}\) and \(\mathfrak{c} / \mathfrak{a}\) if and only if \(\mathbf{A} = \mathfrak{b} + \mathfrak{c}\) and \(\mathfrak{b} \cap \mathfrak{c} = \mathfrak{a}\) . As a result:

Lemma 2. --- The module A/a is indecomposable if and only if a ≠ A and there is no pair (b, c) of ideals of A, distinct from A and a, such that A = b + c and b ∩ c = a.

PROPOSITION 7. --- Let A be a commutative ring, let p be a prime ideal of A (I, p.~117, Def. 3), and let q be an ideal of A contained in p.~Suppose that for every \(x \in p\) there exists an integer n \textgreater{} 0 such that \(x^{n} \in q\) . Then the A-module A/q is indecomposable.

Let \(\mathfrak{b}\) and \(\mathfrak{c}\) be two ideals of \(\mathbf{A}\) , such that \(\mathbf{A} = \mathfrak{b} + \mathfrak{c}\) and \(\mathfrak{b} \cap \mathfrak{c} = \mathfrak{q}\) . Then \(\mathfrak{bc} \subset \mathfrak{b} \cap \mathfrak{c} = \mathfrak{q} \subset \mathfrak{p}\) ; if \(x \notin \mathfrak{p}\) and \(x \in \mathfrak{c}\) , then \(xb \subset \mathfrak{p}\) , so \(\mathfrak{b} \subset \mathfrak{p}\) (I, p.~116, Prop. 4); hence either \(\mathfrak{b} \subset \mathfrak{p}\) or \(\mathfrak{c} \subset \mathfrak{p}\) . Suppose for example that \(\mathfrak{c} \subset \mathfrak{p}\) , so \(\mathfrak{b} + \mathfrak{p} = \mathbf{A}\) ; then there exist \(x \in \mathfrak{b}\) and \(y \in \mathfrak{p}\) such that \(1 = x + y\) ; let \(n \in \mathbf{N}\) be such that \(y^n \in \mathfrak{q}\) ; then \(1 = (x + y)^n\) , so \(1 \in x\mathbf{A} + y^n\mathbf{A} \subset \mathfrak{b} + \mathfrak{q} \subset \mathfrak{b}\) , so \(\mathfrak{b} = \mathbf{A}\) . Lemma 2 now shows that \(\mathbf{A}/\mathfrak{q}\) is indecomposable.

Now suppose that A is a principal ideal domain; by VII, p.~2, Prop. 2, the prime ideals of A are the ideals \((p)\) , where p is an irreducible element of A, and the ideal 0; by the previous proposition, the modules A and \(\mathbf{A}/(p^{n})\) , for p irreducible and n \textgreater{} 0, are indecomposable. Since every cyclic module is a direct sum of modules of this type (VII, p.~3, Prop. 4) and since every finitely generated A-module is a direct sum of cyclic modules (VII, p.~19, Th. 2), we deduce:

PROPOSITION 8. --- Let A be a principal ideal domain and let M be a finitely generated A-module.

\begin{enumerate}
\def\labelenumi{\alph{enumi})}
\item
  M is indecomposable if and only if it is isomorphic to A or to a module of the form \(\mathbf{A} / (p^n)\) , where \(p\) is an irreducible element of A and \(n > 0\) is an integer.
\item
  M is a direct sum of a finite family of indecomposable submodules.
\end{enumerate}

Part b) of the above proposition can be made more precise as follows :

PROPOSITION 9. --- Let A be a principal ideal domain, let P be a system of representatives of irreducible elements of A and let M be a finitely generated A-module. Then there exist positive integers \(m(0)\) and \(m(p^{n})\) ( \(p \in P\) , n \textgreater{} 0), uniquely determined by M and zero except for finitely many of them, such that M is isomorphic to the direct sum of \(\mathbf{A}^{m(0)}\) and the \((\mathbf{A}/(p^{n}))^{m(p^{n})}\) ( \(p \in P\) , n \textgreater{} 0).

The existence of the integers \(m(0)\) and \(m(p^n)\) ( \(p \in P\) , \(n > 0\) ) follows from Prop. 8. The integer \(m(0)\) is uniquely determined: it is the rank of the free module which is the quotient of \(M\) by its torsion submodule. Finally, the \(p\) -primary component of \(M\) is isomorphic to the direct sum of the \((A/(p^n))^{m(p^n)}\) ; since the family of ideals \((p^n)\) ( \(n \geqslant 1\) ) is totally ordered by inclusion, the uniqueness of the \(m(p^n)\) follows from the Cor. to Prop. 2 of VII, p.~16.

DEFINITION 4. --- In the notation of Prop. 9, those ideals \((p^{n})\) ( \(p \in P\) , \(n \geqslant 1\) an integer) such that \(m(p^{n}) > 0\) are called the elementary divisors of the module M, and the integers \(m(p^{n})\) are called their multiplicities; if the integer \(m(0)\) is \textgreater{} 0, it is called the multiplicity of the elementary divisor 0.

As for the invariant factors (VII, p.~19, Def. 1), when A = Z or A = K{[}X{]} (K a commutative field), then the canonical generator of the ideal \((p^{n})\) (a positive integer or a monic polynomial) is also called, by abuse of language, an elementary divisor of the finitely generated module M.

Remarks. --- 1) If M is a finite abelian group, then its structure can be described by writing down its elementary divisors, each repeated as often as its multiplicity. We will say, for example, that M is « of type (2, 2, 4, 27, 27, 25) » (or that it is « a group (2, 2, 4, 27, 27, 25) ») if it is isomorphic to the product of two groups \(\mathbf{Z}/(2)\) , one group \(\mathbf{Z}/(2^{2})\) , two groups \(\mathbf{Z}/(3^{3})\) and one group \(\mathbf{Z}/(5^{2})\) .

\begin{enumerate}
\def\labelenumi{\arabic{enumi})}
\setcounter{enumi}{1}
\tightlist
\item
  If a finitely generated torsion module \(\mathbf{M}\) over a principal ideal domain \(\mathbf{A}\) is given as a direct sum of cyclic modules isomorphic to \(\mathbf{A} / (a_i)\) (in particular if the invariant factors of \(\mathbf{M}\) are known), then the elementary divisors of \(\mathbf{M}\) , and their multiplicities, can be determined by noticing that \(\mathbf{A} / (a)\) is isomorphic to the product of the \(\mathbf{A} / (p^{n(p)})\) , where \(a = \varepsilon \prod_{p \in \mathbb{P}} p^{n(p)}\) is the decomposition of \(a\) into irreducible
\end{enumerate}

factors (VII, p.~3). Let us study for example the multiplicative group G(464 600), where G(n) denotes the multiplicative group \((\mathbf{Z}/n\mathbf{Z})^{*}\) (VII, p.~12). Since \(464\;600 = 2^{3}\cdot5^{2}\cdot23\cdot101\) , this group is isomorphic to the product of the groups G( \(2^{3}\) ), G( \(5^{2}\) ), G(23) and G(101) (VII, p.~13, Th. 3); now the last three groups are cyclic of orders 20, 22 and 100, and G( \(2^{3}\) ) is the product of two cyclic groups of order 2 (loc. cit.) ; since \(20 = 2^{2}\cdot5\) , 22 = \(2\cdot 11\) and \(100 = 2^{2}\cdot5^{2}\) , the group G(464 600) is of type (2, 2, 2, \(2^{2}\) , \(2^{2}\) , 5, \(5^{2}\) , 11).

\begin{enumerate}
\def\labelenumi{\arabic{enumi})}
\setcounter{enumi}{2}
\tightlist
\item
  To calculate the invariant factors of a torsion module whose elementary divisors are known, we again lean on the fact that, if the \(a_{i}\) are pairwise coprime elements of A, then the product \(\prod_{i}\mathbf{A}/(a_{i})\) is a cyclic module isomorphic to
\end{enumerate}

\(A/(a_{1}a_{2}\ldots a_{n})\) (VII, p.~3, Prop. 4). Let us illustrate the method by looking at the example of the group G(464 600) = M : write the elementary divisors \(p^{n}\) of M which are powers of the same irreducible p on the same line, beginning with those of greatest exponent ; extend these lines to lines of equal length by putting in 1's where necessary :

\[
\begin{array}{c c c c c} 2 ^ {2}, & 2 ^ {2}, & 2, & 2, & 2 \\ 5 ^ {2}, & 5, & 1, & 1, & 1 \\ 11, & 1, & 1, & 1, & 1. \end{array}
\]

Then the invariant factors are the products of elements in the same column : 1100, 20, 2, 2, 2. Indeed M is isomorphic to a product of cyclic groups of orders 1100, 20, 2, 2, 2 by Prop. 4 of VII, p.~3 ; since each of these orders is a multiple of the next, these are the invariant factors of M (VII, p.~22, Th. 3).

\begin{enumerate}
\def\labelenumi{\arabic{enumi})}
\setcounter{enumi}{3}
\tightlist
\item
  An A-module is called simple (I, p.~37) if it is nonzero and has no submodules other than itself and 0; it is then necessarily cyclic, so finitely generated, and indecomposable; since the modules \(\mathbf{A} / (p^n)\) are not simple for \(n \neq 1\) , while the modules \(\mathbf{A} / (p)\) are, and since \(\mathbf{A}\) is simple only if the ring \(\mathbf{A}\) is a field, we deduce that the simple A-modules are:
\end{enumerate}

\begin{enumerate}
\def\labelenumi{\alph{enumi})}
\item
  free modules of rank 1, when \(\mathbf{A}\) is a field;
\item
  modules isomorphic to quotients \(\mathbf{A} / (p)\) , where \(p\) is an irreducible element of \(\mathbf{A}\) , when \(\mathbf{A}\) is not a field.
\end{enumerate}

